% =========================================================
% CONFIGURACIÓN DEL CUERPO PRINCIPAL DEL DOCUMENTO
% =========================================================
% Este fichero contiene únicamente la configuración visual del
% cuerpo principal del trabajo:
% - cabeceras y pies de página;
% - formato de títulos de sección;
% - espaciados de títulos;
% - separación vertical de captions.
%
% No incluye configuración del índice porque esa ya está
% centralizada en main.tex.

% =========================================================
% GEOMETRÍA Y CABECERA
% =========================================================

\restoregeometry
% Restaura la geometría principal definida en main.tex.
% Es útil si en páginas anteriores (portada, preliminares, etc.)
% se ha cambiado temporalmente la geometría.

\setlength{\headsep}{4pt}
% Ajusta la distancia entre la cabecera y el inicio del texto.

% =========================================================
% ENCABEZADO Y PIE DE PÁGINA
% =========================================================

\pagestyle{fancy}
% Activa el estilo de página fancyhdr para el cuerpo del documento.

\fancyhf{}
% Limpia completamente cabeceras y pies anteriores antes de redefinirlos.

\fancyhfoffset[L]{0pt}
% Elimina desplazamiento extra en la cabecera/pie por el lado izquierdo.

\fancyhfoffset[R]{0pt}
% Elimina desplazamiento extra en la cabecera/pie por el lado derecho.

\setlength{\headwidth}{\textwidth}
% Hace que la anchura de la cabecera coincida exactamente con el ancho del texto.

\cfoot{\thepage}
% Coloca el número de página centrado en el pie.

% =========================================================
% FORMATO DE TÍTULOS DE CAPÍTULO / SECCIÓN
% =========================================================

\titleformat{\section}[display]
  {\vspace*{0.5cm}\selectfont\bfseries\color{black}}
  {\fontsize{20pt}{20pt}\selectfont Capítulo \thesection}
  {0pt}
  {\fontsize{24pt}{24pt}\selectfont}
% Define el estilo de \section como un título grande en formato display.
% - Añade un pequeño espacio vertical antes del título.
% - Muestra "Capítulo N" en 20pt.
% - Muestra el título del capítulo en 24pt.
% - Usa negrita y color negro.

\titlespacing{\section}
  {0pt}
  {18pt}
  {10pt}
% Controla los espacios de \section:
% - margen izquierdo = 0pt
% - espacio antes del título = 18pt
% - espacio después del título = 10pt

% =========================================================
% FORMATO DE SUBSECCIONES
% =========================================================

\titleformat{\subsection}
  {\normalfont\fontsize{16pt}{20}\bfseries}
  {\thesubsection}
  {0.3em}
  {}
% Define el estilo visual de \subsection:
% - texto normal, negrita
% - tamaño 16pt
% - muestra la numeración automáticamente
% - separa número y título con 0.3em

\titlespacing\subsection{0pt}{12pt plus 4pt minus 2pt}{0pt plus 2pt minus 2pt}
% Controla el espaciado de \subsection:
% - sin margen izquierdo
% - espacio flexible antes del título
% - espacio flexible después del título

% =========================================================
% FORMATO DE SUBSUBSECCIONES
% =========================================================

\titleformat{\subsubsection}
  {\normalfont\fontsize{14pt}{20}\selectfont\bfseries}
  {\thesubsubsection}
  {0.3em}
  {}
% Define el estilo visual de \subsubsection:
% - texto normal
% - tamaño 14pt
% - en negrita
% - muestra la numeración
% - separa número y título con 0.3em

\titlespacing\subsubsection{0pt}{12pt plus 4pt minus 2pt}{0pt plus 2pt minus 2pt}
% Controla el espaciado de \subsubsection.

% =========================================================
% ESPACIADO DE CAPTIONS
% =========================================================

\setlength{\belowcaptionskip}{0pt}
% Elimina espacio extra debajo de tablas y figuras.

\setlength{\abovecaptionskip}{2pt}
% Deja un pequeño espacio encima de los captions.


% =========================================================
% Formato de índices (TOC, LOF, LOT)
% =========================================================



% --- Secciones (Capítulos) ---
\titlecontents{section}
  [0pt] % Sin sangría
  {\fontspec{Calibri}\addvspace{10pt}\fontsize{11pt}{14pt}}
  {Capítulo \thecontentslabel\;} % Ej: "Capítulo 1"
  {}
  {\titlerule*[0.7pc]{.}\contentspage}

% --- Subsecciones (primera versión, redundante) ---
\titlecontents{subsection}
  [2em]
  {\addvspace{2pt}\normalfont}
  {\thecontentslabel\quad}
  {}
  {\titlerule*[0.7pc]{.}\contentspage}

% --- Subsecciones (segunda versión, la que realmente usabas) ---
\titlecontents{subsection}
  [2em]
  {\addvspace{1pt}\sffamily\fontsize{11pt}{12pt}\selectfont}
  {\thecontentslabel\quad}
  {}
  {\titlerule*[0.8pc]{.}\contentspage}

% --- Subsubsecciones ---
\titlecontents{subsubsection}
  [4.0em]
  {\addvspace{1pt}\normalfont\fontsize{10pt}{12pt}\selectfont}
  {\thecontentslabel\quad}
  {}
  {\titlerule*[0.8pc]{.}\contentspage}


\titlecontents{figure}
  [0pt]
  {\color{black}\fontspec{Calibri}\addvspace{5pt}\fontsize{11pt}{14pt}\color{black}}
  {\color{black}Figura \thecontentslabel\quad}
  {\color{black}}
  {\color{black}\titlerule*[0.7pc]{.}\contentspage}

\titlecontents{table}
  [0pt]
  {\color{black}\fontspec{Calibri}\addvspace{5pt}\fontsize{11pt}{14pt}\color{black}}
  {\color{black}Tabla \thecontentslabel\quad}
  {\color{black}}
  {\color{black}\titlerule*[0.7pc]{.}\contentspage}


\addto\captionsspanish{%
  \renewcommand{\contentsname}{Índice de contenidos}%
  \renewcommand{\refname}{Bibliografía}%
}

% Compatibilidad con plantilla previa
\addto\captionsgalician{\renewcommand{\contentsname}{Contenido}}